
Descriptive statistics confirmed that all three canonical stylized facts were present in both the in-sample (2014--2024, $T=2{,}766$) and out-of-sample (2025, $T=249$) windows (Table~\ref{tab:descriptive_stats}): excess kurtosis of 7.7 (IS) and 6.9 (OoS) with Jarque-Bera rejection ($p<0.001$), persistent volatility clustering (Ljung-Box on $|G_t|$ rejected, $p<0.001$), and near-zero linear autocorrelation in raw returns \cite{fama_efficient_1970}.

% ── Table 1: SPY Descriptive Statistics ──────────────────────────────────────
\begin{table}[tp]
\centering
\caption{Descriptive statistics and stylized-facts tests for SPY daily excess growth rates.
    In-sample: 2014--2024 ($T=2{,}766$); out-of-sample: 2025 ($T=249$).
    Gaussian and Laplace columns show maximum likelihood estimation (MLE) fits to the in-sample data.
    JB = Jarque-Bera normality test.
    LB = Ljung-Box autocorrelation test at lag 20.
    $^\ast$ denotes rejection at $\alpha=0.05$.}
\label{tab:descriptive_stats}
\small
\begin{tabular}{@{}lcccc@{}}
\toprule
\textbf{Statistic} & \textbf{IS Observed} & \textbf{OoS Observed} & \textbf{Gaussian} & \textbf{Laplace} \\ \midrule
Mean (annualized, \%)       & 6.31   & 11.60  & 6.31   & 14.19  \\
Std Dev (annualized, \%)    & 214.50 & 220.62 & 214.46 & 205.93 \\
Skewness                    & $-0.753$ & $-1.093$ & 0    & 0      \\
Excess Kurtosis             & 7.715  & 6.867  & 0      & 3      \\
JB normality test           & reject$^\ast$ ($p<0.001$) & reject$^\ast$ ($p<0.001$) & n/a   & n/a   \\
LB test on $G_t$ (lag 20)   & reject$^\ast$ ($p<0.001$) & reject$^\ast$ ($p=0.019$) & n/a   & n/a   \\
LB test on $|G_t|$ (lag 20) & reject$^\ast$ ($p<0.001$) & reject$^\ast$ ($p<0.001$) & n/a   & n/a   \\ \bottomrule
\end{tabular}
\end{table}


\subsection{Jump hyperparameter optimization and state resolution}
With the stylized facts confirmed in both evaluation windows, we turned to calibrating the jump mechanism. A multi-objective grid search over $\epsilon\in[10^{-4},\,2.5\times10^{-2}]$ and $\lambda\in[10,\,160]$, with 200 simulated paths per grid point, identified the optimal values $\epsilon^{*}=10^{-4}$ and $\lambda^{*}=100$ (Figure~\ref{fig:parameter_sweep}, Online Appendix~S2). The optimal $\lambda^{*}=100$ was an interior point of the $\lambda$ grid, confirming that the mean jump duration was identifiable from data; $\epsilon^{*}=10^{-4}$ fell at the lower boundary of the $\epsilon$ grid, indicating that the data consistently preferred rarer tail-entry events over the range explored. At these optimal values, the ACF of $|G_t|$ closely tracked the observed SPY ACF across lags 1--252, confirming that the jump mechanism partially reproduced the empirical ARCH effect.

The primary analysis used $N=100$ states, but a sensitivity study over $N\in\{30,60,90,100,150,200\}$ confirmed robustness: KS and AD pass rates remained stable for both HMM-NJ and HMM-WJ across this range, with every state receiving adequate statistical support (Table~\ref{tab:sensitivity_results}, Online Appendix~S3). At $N=350$, certain states were never visited in the historical record, establishing a practical upper bound on state resolution.


\subsection{In-sample distributional and temporal quality}
Having established that the model's hyperparameters were identifiable and its state resolution robust, we benchmarked HMM-WJ against seven alternative generators on 1,000 synthetic paths of 2,766 trading days each (Table~\ref{tab:model_comparison}, Figure~\ref{fig:model_comparison}). Bootstrap resampling set the non-parametric ceiling (100\% KS, 100\% coverage), while the Gaussian i.i.d.\ baseline (0\% KS) confirmed that thin-tailed models were categorically inadequate. Between these anchors, every model excelled on one quality dimension but failed on another. GARCH(1,1) achieved the best temporal fidelity (ACF-MAE 0.031) but catastrophic distributional failure (5.5\% KS), indicating a systematic shape mismatch despite its well-known volatility clustering properties. The GRU closed 83\% of the ACF-MAE gap but suffered variance collapse (0.6\% KS). The HSMM occupied a middle ground (82\% KS) but its coarse 8-state partition could not capture the full growth rate distribution, as reflected by a 38\% kurtosis gap and only 68.7\% quantile coverage; empirical dwell times of 1.1--1.8 steps caused the semi-Markov structure to collapse to memoryless behavior. This pattern illustrates the difficulty of simultaneously capturing both distributional and temporal quality with a single paradigm.

Both HMM variants achieved the highest distributional fidelity among all parametric models, with HMM-WJ's pass rates 2--8 percentage points below HMM-NJ's, a cost of the jump mechanism occasionally overweighting tail states. The Student-t($\nu=5$) emissions closely matched the observed kurtosis (within 1.3\% for HMM-WJ), a substantial improvement over Normal emissions which underestimated kurtosis by approximately 29\%. Negative skewness ($-0.75$ observed) was also well reproduced by both HMM variants, arising from the asymmetric state-occupancy distribution inherited from the Laplace partition rather than any explicit skewness parameter; symmetric-innovation models (GARCH, Gaussian, Laplace) all produced near-zero skewness by construction. The continuous distance measures corroborated these findings: HMM-NJ and HMM-WJ stayed closest to the Bootstrap ceiling on both $W_1$ and $H$, while GARCH's $W_1$ deteriorated by 67\% from IS to OoS, exposing a tail misspecification that binary pass rates obscured.

On the temporal dimension, we measured ACF-MAE of $|G_t|$ over lags 1--252 (Figure~\ref{fig:model_comparison}, panel~b). HMM-NJ scored 0.059, closing only 3\% of the gap between the i.i.d.\ floor (0.060) and the GARCH ceiling (0.031), confirming that the standard transition matrix alone could not generate persistent volatility clustering. HMM-WJ scored 0.052, closing 28\% of the gap, the only non-GARCH, non-neural model to reduce ACF-MAE meaningfully below the i.i.d.\ floor. The improvement arose because approximately 24\% of HMM-WJ paths contained at least one Poisson jump event that sustained $\text{ACF}(|G_t|)$ decay across all lags, while the remaining 76\% behaved identically to HMM-NJ. The jump frequency is continuously tunable via $\epsilon$; the grid search selected $\epsilon^{*}=10^{-4}$ because this value minimized the joint objective (Equation~\ref{eq:grid_search}), making the 24\% jump rate the empirically optimal operating point for SPY.


% ── Table 2: Model Comparison ───────────────────────────────────────────────
\begin{table}[tp]
\centering
\caption{In-sample and out-of-sample model comparison for SPY (1{,}000 simulated paths,
    significance level $\alpha=0.05$).
    Bootstrap: i.i.d.\ resample of training data.
    Gaussian/Laplace: i.i.d.\ draws from MLE-fitted distributions.
    GARCH: GARCH(1,1) estimated by quasi-maximum likelihood.
    GRU: 2-layer gated recurrent unit \cite{cho_gru_2014} with Gaussian output head (autoregressive generation).
    HSMM: hidden semi-Markov model with $K=8$ Laplace quantile states,
    negative-binomial dwell times, and Student-t($\nu=5$) emissions \cite{bulla_stylized_2006}.
    HMM-NJ: hidden Markov model ($N=100$ states) without jumps.
    HMM-WJ: hybrid model with Poisson jump-duration extension.
    ACF-MAE: mean absolute error of the ACF of $|G_t|$ over lags 1--252.
    Wasserstein-1: mean absolute difference of sorted empirical quantiles (Eq.~\ref{eq:wasserstein}); units are daily excess growth rates.
    Hellinger: histogram-based distributional overlap distance in $[0,1]$ (Eq.~\ref{eq:hellinger}); $H=0$ is identical, $H=1$ is disjoint.
    Coverage: fraction of 99 empirical quantiles (1st--99th) falling within the
    [5th, 95th] percentile envelope of the corresponding synthetic quantiles.
    Standard errors in parentheses: binomial SE for pass rates,
    $\text{std}/\!\sqrt{n}$ for kurtosis, Wasserstein-1, and Hellinger,
    bootstrap SE ($B=500$) for ACF-MAE and coverage.
    Arrows indicate preferred direction: $\uparrow$~higher is better,
    $\downarrow$~lower is better, $\approx$~closer to observed is better.}
\label{tab:model_comparison}
\resizebox{\columnwidth}{!}{%
\scriptsize
\begin{tabular}{@{}lcccccccc@{}}
\toprule
\textbf{Metric} & \textbf{Bootstrap} & \textbf{Gaussian} & \textbf{Laplace} & \textbf{GARCH(1,1)} & \textbf{GRU} & \textbf{HSMM} & \textbf{HMM-NJ} & \textbf{HMM-WJ} \\ \midrule
\multicolumn{9}{l}{\textit{In-sample: 2{,}766 trading days (2014--2024)}} \\[2pt]
KS pass rate (\%) $\uparrow$            & 100.0 ({$<$}0.1) &   0.0 ({$<$}0.1) &  44.0 (1.6) &   5.5 (0.7) &   0.6 (0.2) &  82.0 (1.2) &  \best{99.7} (0.2) &  97.6 (0.5) \\
AD pass rate (\%) $\uparrow$            &  99.7 (0.2)      &   0.0 ({$<$}0.1) &  43.3 (1.6) &   1.9 (0.4) &   0.2 (0.1) &  42.5 (1.6) &  \best{99.1} (0.3) &  91.3 (0.9) \\
Excess kurtosis (observed)   & \multicolumn{8}{c}{7.715} \\
Excess kurtosis (simulated) $\approx$  &   7.6 (0.08) & $-$0.0 ({$<$}0.01) &   3.0 (0.02) &   \best{8.2} (0.37) &   5.7 (0.16) &   4.8 (0.08) &   8.1 (0.14) &   7.6 (0.14) \\
ACF-MAE $\downarrow$                      & 0.060 ({$<$}0.001) & 0.060 ({$<$}0.001) & 0.060 ({$<$}0.001) & \best{0.031} (0.001) & 0.036 ({$<$}0.001) & 0.059 ({$<$}0.001) & 0.059 ({$<$}0.001) & 0.052 ({$<$}0.001) \\
Coverage (\%) $\uparrow$                & 100.0 ({$<$}0.1)   &  13.1 (0.1)       &  37.4 (1.5) &  29.3 (1.4) &  17.2 (0.5) &  68.7 (1.1) & \best{100.0} ({$<$}0.1)   & \best{100.0} ({$<$}0.1) \\
Wasserstein-1 $\downarrow$              & 0.062 (0.001) & 0.399 (0.001) & 0.138 (0.001) & 0.304 (0.006) & 0.421 (0.003) & 0.176 (0.001) & \best{0.081} (0.001) & 0.101 (0.001) \\
Hellinger dist $\downarrow$             & 0.042 ({$<$}0.001) & 0.148 ({$<$}0.001) & \best{0.072} ({$<$}0.001) & 0.100 (0.001) & 0.134 (0.001) & 0.113 ({$<$}0.001) & 0.073 ({$<$}0.001) & 0.075 ({$<$}0.001) \\[4pt]
\multicolumn{9}{l}{\textit{Out-of-sample: 249 trading days (2025)}} \\[2pt]
KS pass rate (\%) $\uparrow$            &  97.4 (0.5) &  62.2 (1.5) &  88.0 (1.0) &  80.3 (1.3) &  71.8 (1.4) &  96.2 (0.6) &  \best{96.7} (0.6) &  94.4 (0.7) \\
AD pass rate (\%) $\uparrow$            &  98.5 (0.4) &  46.3 (1.6) &  92.9 (0.8) &  72.0 (1.4) &  44.8 (1.6) &  \best{96.7} (0.6) &  \best{96.7} (0.6) &  95.1 (0.7) \\
Excess kurtosis (observed)   & \multicolumn{8}{c}{6.867} \\
Excess kurtosis (simulated) $\approx$  &   6.0 (0.18) & $-$0.0 ({$<$}0.01) &   2.7 (0.05) &   1.6 (0.07) &   2.9 (0.10) &   4.1 (0.13) &   \best{6.4} (0.22) &   6.1 (0.13) \\
ACF-MAE $\downarrow$                      & 0.043 ({$<$}0.001) & 0.043 ({$<$}0.001) & 0.043 ({$<$}0.001) & \best{0.026} ({$<$}0.001) & 0.033 ({$<$}0.001) & 0.042 ({$<$}0.001) & 0.041 ({$<$}0.001) & 0.039 ({$<$}0.001) \\
Coverage (\%) $\uparrow$                & 100.0 ({$<$}0.1) &  36.4 (1.9) &  74.7 (1.0) &  96.0 (1.3) &  77.8 (2.9) & \best{100.0} (0.2) & \best{100.0} ({$<$}0.1) & \best{100.0} ({$<$}0.1) \\
Wasserstein-1 $\downarrow$              & 0.232 (0.002) & 0.452 (0.002) & 0.263 (0.002) & 0.507 (0.018) & 0.488 (0.005) & 0.287 (0.002) & \best{0.258} (0.002) & 0.282 (0.006) \\
Hellinger dist $\downarrow$             & 0.205 (0.001) & 0.235 (0.001) & 0.211 (0.001) & 0.232 (0.001) & 0.240 (0.001) & 0.239 (0.001) & \best{0.207} (0.001) & 0.210 (0.001) \\[4pt]
\midrule
Parameters estimated\tablefootnote{The $N\times N$ transition matrix in HMM-NJ and HMM-WJ is computed by direct counting of observed state transitions; it is a sufficient statistic of the data, not a fitted parameter. The counts listed here reflect only the scalar parameters requiring estimation (MLE or grid search).}         &     0 &     2 &     2 &     3 & 37{,}954 &    18 &     2 &     4 \\ \bottomrule
\end{tabular}}
\end{table}


% ── Figure 4: Head-to-Head Model Comparison ──────────────────────────────────
\begin{figure}[tp]
    \centering
    \includegraphics[width=\textwidth]{Fig4-Model-Comparison.pdf}
    \Description{Three-panel figure comparing generative models for SPY. Panel (a) shows marginal density histograms of excess growth rates with KS pass rates annotated; HMM variants closely match the empirical heavy-tailed distribution while GARCH and Gaussian fail. Panel (b) shows autocorrelation functions of absolute returns at lags 1 to 252; HMM-WJ reproduces slow ACF decay while HMM-NJ and i.i.d. baselines collapse near zero. Panel (c) shows tail Q-Q plots at the 0.1st through 99.9th percentile region; HMM-WJ provides the closest quantile match in the extremes.}
    \caption{Head-to-head in-sample model comparison for SPY ($N=100$, 1{,}000 simulated paths).
    Panel~(a): marginal density of excess growth rates with IS KS pass rates annotated.
    GARCH(1,1) fails the two-sample KS test on 95\% of paths (pass rate 5.5\%), indicating a
    systematic shape mismatch with the empirical distribution.  Both HMM variants pass at
    rates $\geq\!97\%$, confirming that the Student-t emission structure adequately captures the observed
    heavy tails.  Out-of-sample, HMM-WJ maintains 94\% KS pass rate, though HMM-NJ
    achieves a higher 97\%; GARCH recovers partially (OoS 80\%).
    Panel~(b): autocorrelation function of absolute excess growth rates,
    $\mathrm{ACF}(|G_t|)$, at lags 1--252 (shaded bands: 10th--90th percentile across paths).
    HMM-NJ is structurally incapable of producing persistent volatility clustering:
    without a jump mechanism the latent states are i.i.d.\ conditional on the Markov chain,
    so $\mathrm{ACF}(|G_t|)\approx 0$ for all lags beyond one (dotted curve).
    HMM-WJ generates a \emph{mixture} of two path families under the same model:
    approximately 76\% of IS paths contain no Poisson jump and therefore behave
    identically to HMM-NJ (dashed curve, near zero); the remaining $\sim$24\% of paths
    contain at least one jump event and sustain substantial $\mathrm{ACF}(|G_t|)$ decay
    across all lags (solid navy curve $\pm$ band).
    The jump frequency, and hence the fraction of the ensemble exhibiting slow ACF decay, is
    controlled by the tail-entry probability $\epsilon$, making the volatility-clustering strength
    directly tunable.
    Panel~(c): tail Q-Q plot at the 0.1st--99.9th percentile region; HMM-WJ provides the
    closest mean quantile match in the extreme tails.}
    \label{fig:model_comparison}
\end{figure}

\subsection{Out-of-sample evaluation and model tradeoffs}
The in-sample results demonstrated that HMM-WJ was the only model to improve on both distributional and temporal dimensions; the out-of-sample window tested whether this advantage survived on unseen data. We evaluated all eight generators on a held-out window of 249 trading days (full calendar year 2025) using the same metrics (Table~\ref{tab:model_comparison}, Figure~\ref{fig:statistical_validation}). No single model dominated all quality dimensions simultaneously. GARCH(1,1) achieved the best temporal fidelity (closing 100\% of the ACF-MAE gap) but the worst distributional fit among parametric models (5.5\% of the Bootstrap KS ceiling). The GRU closed 83\% of the ACF gap but achieved the worst distributional fit overall (0.6\% of the KS ceiling), demonstrating that even 38{,}000 trainable parameters could not jointly capture both quality dimensions. HMM-NJ achieved the highest distributional pass rates (99.7\% of the KS ceiling) but could not reproduce volatility clustering (closing only 3\% of the ACF gap). HMM-WJ occupied the Pareto frontier: it closed 28\% of the ACF gap while maintaining distributional pass rates above 91\% of the ceiling, the only model to improve meaningfully on both dimensions simultaneously.

Out-of-sample, both HMM variants generalized robustly: HMM-NJ achieved 96.7\% KS and HMM-WJ achieved 94.4\% KS, with simulated kurtosis of 6.4 and 6.1 against an observed 6.9 (Table~\ref{tab:model_comparison}, OoS panel). GARCH's $W_1$ deteriorated by 67\% from IS to OoS, exposing a tail misspecification that binary pass rates obscured. The OoS ordering reversed the IS pattern (HMM-NJ now outperformed HMM-WJ), suggesting mild overfitting of the IS-calibrated jump parameters. Cross-asset generalization to NVDA, JNJ, and JPM confirmed that the jump mechanism automatically adapted its duration $\lambda^{*}$ to each asset's clustering intensity (Online Appendix~S5, Table~\ref{tab:cross_asset}).

% ── Figure 6: Statistical Validation ─────────────────────────────────────────
\begin{figure}[tp]
    \centering
    \includegraphics[width=\textwidth]{Fig6-Statistical-Validation.pdf}
    \Description{Four-panel statistical validation figure for HMM-WJ. Panel (a) shows in-sample KS p-value histograms for jump-containing paths. Panel (b) shows out-of-sample KS p-value histograms for all paths. Panel (c) shows the out-of-sample marginal density fan chart with the simulation envelope. Panel (d) shows the in-sample ACF of absolute returns for jump-containing paths, demonstrating that the jump mechanism reproduces volatility clustering.}
    \caption{Statistical validation of HMM-WJ ($N=100$, $\alpha=0.05$).
    Panels~(a) and~(d) condition on the $\sim$24\% of IS paths containing jump events.
    (a)~IS KS $p$-values for 239 jump paths (pass rate 93.7\%).
    (b)~OoS KS $p$-values for all 1{,}000 paths (pass rate 95.4\%).
    (c)~OoS density fan chart; observed density falls within the simulation envelope.
    (d)~IS ACF of $|G_t|$ for jump paths; the mean closely tracks the observed ACF.
    The all-paths vs.\ jump-paths ACF comparison is in Online Appendix~S7.}
    \label{fig:statistical_validation}
\end{figure}



\subsection{Multi-asset extension}
The single-asset results confirmed distributional and temporal fidelity for SPY; the remaining question was whether the framework could scale to a full asset universe. We propagated HMM-WJ SPY paths through a Single-Index Model (SIM) for 424 S\&P~500 constituents, with idiosyncratic shocks drawn by resampling empirical residuals. The SIM yielded a median IS KS pass rate of 66.7\% and a mean of 58.4\% across the asset universe (well below the 95\% expected under a correct model); high-$\beta$ assets closely tracking SPY retained pass rates above 80\%, while degradation concentrated among low-$R^2$ assets, confirming that the HMM generative engine was not the bottleneck. Out-of-sample, the mean KS pass rate reached 82.1\% across 417 surviving assets. Full cross-sectional results, sector-level analysis, and the SIM $R^2$ distribution are reported in Online Appendix~S8 (Figures~\ref{fig:multi_asset} and~\ref{fig:multi_asset_oos}, Table~\ref{tab:multi_asset}). To determine whether richer dependence specifications could recover the fidelity lost by the SIM's linear assumption, we compared four dependence models on a four-asset portfolio (SPY, NVDA, JNJ, JPM): (i)~a Gaussian copula, (ii)~a Student-t copula, (iii)~a C-vine copula with AIC-selected bivariate families, and (iv)~the Single-Index Model with bootstrap residuals (Figure~\ref{fig:price_trajectories}). We fitted an independent HMM to each asset, fitted each dependence model to the in-sample growth rate matrix, generated 1,000 synthetic paths of 2,766 steps, and evaluated both marginal fidelity (per-asset KS pass rates) and cross-asset dependence quality (correlation reproduction). For the copula models, simulation proceeded by rank reordering: each asset's independently simulated HMM values were rearranged in time so that the cross-asset rank structure matched the fitted copula, preserving each marginal distribution exactly while injecting the desired correlation and tail dependence (Online Appendix~S6).


% ── Table 5: Copula Comparison ────────────────────────────────────────────
\begin{table}[tp]
\centering
\caption{Dependence model comparison for a four-asset portfolio (SPY, NVDA, JNJ, JPM).
    Each model uses per-asset HMM marginals ($N=100$, Student-t emissions with $\nu=5$);
    only the dependence structure differs.
    1{,}000 simulated paths, $\alpha=0.05$.
    Corr Frobenius: root-sum-squared error of upper-triangle correlation entries
    (observed vs.\ mean simulated). Pairwise corr MAE: mean absolute error across
    all six asset pairs. SIM results exclude SPY (market factor).}
\label{tab:copula_comparison}
\small
\begin{tabular}{@{}lcccc@{}}
\toprule
 & \textbf{SIM} & \textbf{Gaussian} & \textbf{Student-t} & \textbf{Vine} \\ \midrule
\multicolumn{5}{l}{\textit{Marginal KS pass rate (\%)}} \\[2pt]
SPY   & n/a           & 98.2 (0.4) & \textbf{98.9} (0.3) & 96.3 (0.6) \\
NVDA  & 79.5 (1.3)    & 96.6 (0.6) & \textbf{99.1} (0.3) & 98.6 (0.4) \\
JNJ   & 92.1 (0.9)    & \textbf{98.3} (0.4) & 97.3 (0.5) & 99.2 (0.3) \\
JPM   & 27.5 (1.4)    & \textbf{97.8} (0.5) & 95.1 (0.7) & 97.4 (0.5) \\[4pt]
\multicolumn{5}{l}{\textit{Cross-asset dependence quality}} \\[2pt]
Corr Frobenius error $\downarrow$  & n/a  & 0.168 (0.001) & \textbf{0.134} (0.001) & 0.287 (0.002) \\
Pairwise corr MAE $\downarrow$     & n/a  & 0.055         & \textbf{0.039}         & 0.099         \\
\bottomrule
\end{tabular}
\end{table}

We observed a clear separation between copula-based and factor-based approaches on marginal fidelity (Table~\ref{tab:copula_comparison}, Figure~\ref{fig:copula_comparison}, panel~a). All three copula specifications achieved KS pass rates above 95\% for every ticker, which confirmed that the rank-reordering mechanism preserved each asset's HMM marginal distribution as expected. The SIM, by contrast, produced KS pass rates of 79.5\% for NVDA, 92.1\% for JNJ, and only 27.5\% for JPM, which reflected the inability of a linear factor model to reproduce asset-specific tail structure when the factor loading was moderate ($R^2 = 0.48$ for JPM). This gap demonstrated that the multi-asset quality limitation identified in the 424-asset SIM analysis was attributable to the dependence specification, not to the underlying HMM marginal models. While copulas preserved marginals, the key question was whether they also captured cross-asset dependence. We measured this using the Frobenius norm of the correlation error and the mean absolute error across all six pairwise correlations (Figure~\ref{fig:copula_comparison}, panel~b). The Student-t copula achieved the best correlation reproduction (Frobenius error 0.134, SE~0.001; pairwise MAE 0.039), followed by the Gaussian copula (0.168; 0.055) and the vine copula (0.287; 0.099). The vine copula's poor correlation reproduction in this four-asset setting was attributable to overfitting: with only six bivariate edges and five candidate families per edge, the AIC-selected model captured local pairwise structure at the expense of global correlation coherence. The Student-t copula's advantage over the Gaussian copula arose from its ability to model tail dependence; at the fitted $\nu = 5$ degrees of freedom, extreme co-movements across assets occurred substantially more frequently than a Gaussian dependence structure would produce, which better matched the observed concentration of joint drawdowns in the historical data.

% ── Figure: Price Trajectories ─────────────────────────────────────────────
\begin{figure*}[h!]
    \centering
    \includegraphics[width=\textwidth]{Fig-Price-Trajectories.pdf}
    \Description{Three-panel figure showing simulated price trajectories for SPY, NVDA, JNJ, JPM under three dependence specifications.}
    \caption{Simulated price trajectories under three dependence specifications (all starting at \$100).
    (a)~Independent HMM marginals: drawdowns and rallies occur at unrelated times.
    (b)~Single-Index Model: partial co-movement via the market factor.
    (c)~Student-t copula: crash episodes synchronize because the copula reorders each asset's
    independently simulated values so that extreme days coincide.
    Panels~(a) and~(c) use the same growth rate values; only the temporal ordering differs.}
    \label{fig:price_trajectories}
\end{figure*}

The vine copula's underperformance had a revealing explanation: the AIC-based selection procedure chose the Student-t bivariate copula at all six edges across all three vine tree levels ($\Delta\text{AIC} > 60$ over every alternative \cite{burnham_model_2002}), with no Clayton, Gumbel, or Frank families selected. This was unexpected given that the financial risk literature frequently posits asymmetric tail dependence in equity returns \cite{ang_asymmetric_2002, longin_extreme_2001}. The fitted degrees-of-freedom parameter varied systematically: $\nu=4$--$5$ for direct market pairs (SPY-NVDA, SPY-JNJ, SPY-JPM), indicating heavy symmetric tail dependence, and $\nu=10$--$30$ for conditional pairs after removing the market factor, indicating that most tail dependence was attributable to the common market component. Because every vine edge selected the same family, the vine's flexibility provided no advantage over a single Student-t copula while adding estimation noise from the sequential h-function transformations. For these four U.S.\ large-cap equities over this sample period, a single Student-t copula was a parsimonious and sufficient specification; whether this symmetry persists across broader asset universes or crisis periods remains an open question.